% Bab V Implementasi Arsitektur Platform
% Template: TEMPLATE_BAB.md
% Gaya: WRITING_GUIDE.md
\cleardoublepage
\chapter{IMPLEMENTASI ARSITEKTUR PLATFORM DATAOPS DAN MLOPS}
\label{chap:implementasi}

Bab ini membahas implementasi platform DataOps dan MLOps terintegrasi sesuai rancangan pada Bab~\ref{chap:perancangan}. Pembahasan dimulai dari lingkungan implementasi yang menjadi dasar \textit{deployment} komponen, kemudian dilanjutkan dengan implementasi empat sub-sistem mengikuti urutan pemenuhan tujuan T-1 sampai T-4. Bab ini ditutup dengan verifikasi \textit{execution flow} menggunakan satu \textit{use case} sebagai kasus pengujian. Penjabaran detail biner dan konfigurasi disimpan pada repositori platform, sehingga bab ini menyajikan pokok-pokok implementasi yang menjelaskan bentuk akhir rancangan.

\section{Lingkungan Implementasi}
\label{sec:lingkungan}

Lingkungan implementasi platform dibangun pada satu \textit{cluster} Kubernetes \autocite{kubernetes2024} berbasis k3s yang berjalan pada satu \textit{node}. Pilihan ini memenuhi kebutuhan eksekusi pada lingkungan terbatas tanpa mengubah pola \textit{deployment} yang dapat dipakai pada \textit{cluster} produksi dengan banyak \textit{node}. Distribusi k3s membawa \textit{control plane} ringan, \textit{containerd} sebagai \textit{runtime}, dan beberapa \textit{add-on} \textit{built-in} yang dapat dinonaktifkan agar tidak berbenturan dengan komponen pilihan platform. Konfigurasi \textit{node} memakai satu \textit{disk} lokal dengan \textit{ext4} sebagai \textit{filesystem}, dan tata letak \textit{namespace} pada platform dibakukan agar perpindahan ke \textit{cluster} produksi dengan banyak \textit{node} tinggal mengubah berkas \textit{overlay} \textit{Kustomize} tanpa mengganti manifes dasar.

Kode dan konfigurasi proyek ditata pada beberapa direktori utama dengan batas tanggung jawab yang jelas. Tabel~\ref{tab:struktur-direktori} merangkum direktori tersebut beserta fungsinya, dengan \texttt{platform/} sebagai inti yang bersifat domain-agnostik dan direktori \textit{use case} sebagai kasus pengujian yang terpisah. Pemisahan ini memastikan platform dapat dipakai ulang lintas domain tanpa menyentuh berkas yang spesifik pada satu \textit{use case}. Direktori \texttt{experiments/} dan \texttt{materials/} berperan sebagai penunjang verifikasi dan dokumentasi yang tidak menjadi bagian dari \textit{deployment flow} platform.

\begin{table}[!htb]
\centering
\caption{Struktur Direktori Utama Proyek dan Fungsinya}
\label{tab:struktur-direktori}
\footnotesize
\setlength{\tabcolsep}{3pt}
\begin{tabular}{|>{\raggedright\arraybackslash}m{3.1cm}|>{\raggedright\arraybackslash}m{10.3cm}|}
\hline
\textbf{Direktori} & \textbf{Fungsi} \\
\hline
\texttt{platform/} & Kode, manifes, dan konfigurasi seluruh komponen platform DataOps dan MLOps yang bersifat domain-agnostik \\
\hline
\texttt{use-case-crypto/} & \textit{Use case} uji berupa pasar kripto, memuat layanan ingestasi, definisi fitur, \textit{pipeline}, dan \textit{overlay} \textit{Kustomize} yang spesifik domain \\
\hline
\texttt{experiments/} & Skenario verifikasi dan \textit{oracle} pengujian, antara lain uji paritas \textit{batch}-\textit{stream}, beban, \textit{drift}, dan \textit{chaos} \\
\hline
\texttt{materials/} & Dokumen tugas akhir, proposal, referensi, dan berkas laporan \\
\hline
\end{tabular}
\end{table}


\subsection{Batasan Implementasi}
\label{subsec:batasan-implementasi}

Perwujudan platform pada bab ini berjalan di dalam dua batasan implementasi yang dibedakan dari batasan penelitian pada Bab~\ref{chap:pendahuluan}. Batasan implementasi menyangkut pilihan perwujudan yang dapat diganti tanpa mengubah sifat kontribusi selama fungsi yang sama tetap terpenuhi, karena keluaran utama penelitian adalah arsitektur yang dapat dipakai lintas \textit{use case}, bukan implementasi untuk satu lingkungan tertentu. Lingkungan k3s satu \textit{node} dan \textit{version pinning} April 2026 ditetapkan oleh kedua batasan itu sebagai kondisi pengujian yang dapat direproduksi. Selanjutnya, kedua butir itu dirujuk dengan singkatan BI-1 dan BI-2, terutama pada pembahasan hasil di Bab~\ref{chap:evaluasi}.

\vspace{0.5em}
\begin{enumerate}
    \item Platform berjalan di atas k3s pada satu \textit{node} sebagai lingkungan pengembangan dan pengujian, dengan skala diatur oleh HPA, VPA, dan KEDA ketika beban nyata muncul, sedangkan \textit{overlay} Kustomize untuk \textit{environment} \textit{multi-node} tetap dapat ditambahkan tanpa modifikasi basis kode sehingga jumlah \textit{node} bersifat perwujudan dan bukan bagian dari kontribusi.

    \item Versi pustaka dan basis \textit{image container} dipatok pada rilis stabil tertinggi per April 2026 agar pengujian dapat direproduksi pada lingkungan yang sama, dan \textit{version pinning} maupun pilihan serta sintaks komponen tertentu dapat diganti dengan komponen lain berkapabilitas setara tanpa mengubah rancangan.
\end{enumerate}

Batasan implementasi tersebut menjadi konteks pembacaan seluruh subbab implementasi pada bab ini. Penggantian jumlah \textit{node}, versi komponen, atau sintaks \textit{tools} dengan padanan berkapabilitas setara tidak mengubah rancangan pada Bab~\ref{chap:perancangan}. Batasan BI-1 dipetakan kembali pada pelaporan hasil evaluasi lingkungan \textit{node} tunggal, sedangkan BI-2 menjadi acuan reproduksi lingkungan pengujian pada Bab~\ref{chap:evaluasi}. Pemisahan dari batasan penelitian menegaskan bahwa keterbatasan lingkungan pengujian tidak membatasi keberlakuan arsitektur yang dirancang.

\subsection{Konfigurasi Dasar \textit{Cluster}}
\label{subsec:cluster}

\textit{Cluster} menjalankan komponen dasar berikut. Istio \autocite{istio2024} berperan sebagai \textit{service mesh} dengan mTLS dalam mode \textit{strict} pada trafik internal, APISIX \autocite{apisix2025} sebagai \textit{API gateway} pada \textit{mesh edge}, kontrol akses jaringan menggunakan \textit{NetworkPolicy} pada level Kubernetes, dan Kyverno \autocite{kyverno2024} sebagai \textit{policy engine} pada \textit{admission layer}. Falco \autocite{falco2025} memantau perilaku \textit{runtime} pada level \textit{kernel}, sedangkan Trivy Operator \autocite{trivy2025} memindai \textit{image} dan \textit{workload} secara berkala. \textit{Storage class} \textit{default} menggunakan \textit{local-path provisioner} untuk \textit{persistent volume} berbasis penyimpanan lokal, dengan Velero \autocite{velero2025} melakukan pencadangan terjadwal ke MinIO. Chaos Mesh \autocite{chaosmesh2025} disiapkan untuk uji ketahanan dengan eksperimen \textit{fault injection} yang terbatas pada \textit{namespace} tertentu.

Konfigurasi sumber daya memberi prioritas pada satu replika \textit{idle} per \textit{workload}, sementara \textit{Horizontal Pod Autoscaler}, \textit{Vertical Pod Autoscaler} dalam mode rekomendasi, dan KEDA \textit{ScaledObject} disiapkan pada \textit{template} agar \textit{scaling} dapat aktif ketika beban nyata datang. \textit{Resource limits} pada setiap kontainer ditetapkan sesuai panduan Kyverno \texttt{require-resource-limits} untuk mencegah kelebihan kuota \textit{node}. Penjadwalan kelas prioritas memakai \texttt{PriorityClass} platform agar komponen kritis tidak di-\textit{preempt} oleh \textit{workload} sekali pakai. \textit{External Snapshotter} dipasang sebagai prasyarat untuk \textit{snapshot} CSI yang dipakai oleh Velero, sehingga rekonstruksi \textit{cluster} dari pencadangan dapat dilakukan tanpa intervensi manual.

\subsection{Manajemen \textit{Secret} dan Identitas}
\label{subsec:rahasia-identitas}

OpenBao \autocite{openbao2025} bertindak sebagai penyimpan \textit{secret}. \textit{External Secrets Operator} bertugas menyalurkan \textit{secret} dari OpenBao ke \textit{namespace} target tanpa menyalin nilai \textit{secret} ke dalam berkas manifes. Identitas pengguna dikelola oleh dex \autocite{dex2025} sebagai \textit{identity provider} OIDC, dengan oauth2-proxy \autocite{oauth2proxy2025} sebagai \textit{reverse proxy} yang menegakkan autentikasi pada antarmuka aplikasi seperti DataHub, MLflow, Grafana, Argo CD, dan Kubeflow Pipelines. Selain identitas pengguna, kebijakan otorisasi \textit{workload} pada \textit{admission layer} ditegakkan oleh Kyverno \autocite{kyverno2024} sebagai bagian dari dasar keamanan \textit{cluster}.

OpenBao diinisialisasi oleh \textit{Job} \texttt{openbao-bootstrap} yang melakukan \textit{init} dan \textit{auto-unseal} dengan kunci yang tersimpan di MinIO \autocite{minio2024} dengan enkripsi sisi \textit{server}. \textit{External Secrets Operator} dipasang pada \textit{namespace} \texttt{security} dengan \texttt{ClusterSecretStore} tunggal yang berisi konfigurasi koneksi ke OpenBao. SpiceDB \autocite{spicedb2025} dipasang sebagai layanan otorisasi berbasis Zanzibar yang menetapkan hubungan akses pada level objek tata kelola, dengan basis data PostgreSQL melalui CNPG sebagai \textit{datastore}. \textit{Key Encryption Service} dipasang sebagai \textit{sidecar} MinIO untuk enkripsi \textit{secret}, dengan rotasi kunci dijalankan oleh \textit{CronJob} terjadwal.

\subsection{Pola GitOps}
\label{subsec:gitops}

Argo CD \autocite{argocd2024} menjadi satu-satunya \textit{change flow} pada \textit{cluster}. Repositori Git pada Gitea \autocite{gitea2025} bertindak sebagai sumber rujukan konfigurasi. Setiap perubahan dilakukan melalui \textit{commit} pada repositori dan diterapkan ke \textit{cluster} oleh Argo CD melalui rekonsiliasi. Pola ini menjamin reproduksibilitas, menghasilkan jejak audit untuk setiap perubahan, dan memenuhi KF-12 tentang konfigurasi deklaratif. \textit{Pipeline} CI yang dijalankan oleh Tekton \autocite{tekton2024} melakukan \textit{build} \textit{image} dan menjalankan pengujian sebelum \textit{commit} menjadi sumber \textit{deployment}.

\textit{ApplicationSet} Argo CD dipakai untuk mendeklarasikan banyak \texttt{Application} sekaligus dengan satu \textit{template} per kelompok, sehingga penambahan komponen baru tinggal menambahkan \textit{generator} tanpa menambah berkas \texttt{Application} satu per satu. Argo CD diatur dengan mode \texttt{manual-sync} pada \textit{ApplicationSet} \textit{template} untuk mencegah \texttt{auto-prune} yang dapat menghapus sumber daya secara tidak sengaja. Tekton menjalankan \textit{Pipeline} yang menghasilkan \textit{image} ke registri \texttt{localhost:5000} yang berfungsi sebagai \textit{mirror} lokal, dan menulis kembali \texttt{commit} dengan referensi \textit{image} yang sudah diperbarui pada \texttt{values.yaml} \textit{kustomization}. \textit{Pre-apply} dan \textit{post-apply} \textit{hook} pada \texttt{apply-component.sh} menutup celah \textit{race condition} antara CRD dan CR, sehingga seluruh komponen \textit{cluster} dapat dipasang ulang dengan satu perintah Makefile.

\section{Implementasi Arsitektur Terintegrasi}
\label{sec:impl-arsitektur}

Bagian ini mengimplementasikan rancangan pada Bagian~\ref{sec:arsitektur-terintegrasi} sebagai pemenuhan T-1. Komponen disusun ke dalam \textit{namespace} per \textit{layer} dan dihubungkan melalui kontrak API yang stabil. Setiap \textit{namespace} memuat satu kelompok komponen fungsional sehingga kebijakan \textit{NetworkPolicy} dapat ditegakkan dengan model \textit{default-deny} dan pengecualian eksplisit per arah komunikasi. \textit{Deployment} komponen mengikuti pola GitOps melalui Argo CD, sedangkan eksekusi \textit{pipeline} CI mengikuti pola Tekton, dengan setiap komponen tunduk pada \textit{template} \textit{ApplicationSet} agar penambahan komponen baru tidak menambah \textit{boilerplate}. Dari sembilan \textit{layer} rancangan, tiga \textit{layer} operasional telah terimplementasi pada Subbab~\ref{sec:lingkungan}, yaitu \textit{Orchestration and Infrastructure Layer} pada konfigurasi dasar \textit{cluster}, \textit{Platform Security Layer} pada manajemen \textit{secret} dan identitas, serta \textit{GitOps and Continuous Delivery Layer} pada pola GitOps. \textit{Feature Service Layer} dan \textit{Data Governance Layer} diimplementasikan sebagai \textit{sub-sistem} pada Subbab~\ref{sec:impl-feast} dan Subbab~\ref{sec:impl-tata-kelola}, sehingga subbab ini merinci empat \textit{layer} sisanya. Tabel~\ref{tab:impl-layer-map} merangkum bentuk akhir implementasi tersebut dengan memetakan tiap \textit{layer} ke komponen yang terpasang beserta \textit{namespace} tempatnya berjalan, sehingga pembaca dapat memeriksa kelengkapan sembilan \textit{layer} tanpa menelusuri manifes satu per satu.

% Tabel pemetaan implementasi V.2: sembilan layer -> komponen terpasang ->
% namespace. Nama namespace diverifikasi dari platform/components/*
% (kustomization + manifes). Gaya mengikuti struktur_direktori.tex
% (border penuh, footnotesize, tabcolsep 3pt, lebar total ~ \textwidth).
\begin{table}[!htb]
\centering
\caption{Pemetaan \textit{Layer} Platform ke Komponen Terpasang dan \textit{Namespace}}
\label{tab:impl-layer-map}
\footnotesize
\setlength{\tabcolsep}{3pt}
\begin{tabular}{|>{\raggedright\arraybackslash}m{3.4cm}|>{\raggedright\arraybackslash}m{5.9cm}|>{\raggedright\arraybackslash}m{4.0cm}|}
\hline
\centering\arraybackslash\textbf{\textit{Layer}} & \centering\arraybackslash\textbf{Komponen Terpasang} & \centering\arraybackslash\textbf{\textit{Namespace}} \\
\hline
\textit{Orchestration and Infrastructure} & k3s, metrics-server, HPA, VPA, KEDA, Kueue, cert-manager & \texttt{kube-system}, \texttt{common}, \texttt{cert-manager} \\
\hline
\textit{Data Ingestion and Processing} & Kafka (Strimzi), Kafka Connect, Debezium, Karapace, Flink, Spark, dbt, Airflow & \texttt{data-ingestion}, \texttt{data-processing} \\
\hline
\textit{Data Storage} & MinIO, Iceberg, Lakekeeper, lakeFS, ClickHouse, Trino, PostgreSQL (CNPG), MySQL & \texttt{storage} \\
\hline
\textit{Feature Service} & Feast, Valkey, Qdrant & \texttt{model-lifecycle}, \texttt{storage} \\
\hline
\textit{Model Lifecycle} & MLflow, Kubeflow Pipelines, Katib, Kubeflow Trainer, Argo Workflows, KServe, Knative, Argo Rollouts & \texttt{model-lifecycle}, \texttt{model-serving}, \texttt{knative-serving} \\
\hline
\textit{Data Governance} & DataHub, OpenLineage, Great Expectations, OpenSearch & \texttt{data-governance} \\
\hline
\textit{Observability} & Prometheus, Alertmanager, Loki, Tempo, Grafana, OpenTelemetry, Superset, Evidently & \texttt{observability} \\
\hline
\textit{GitOps and Continuous Delivery} & Gitea, Tekton, Argo CD, registri \textit{image} lokal & \texttt{gitops}, \texttt{tekton-pipelines}, \texttt{platform-registry} \\
\hline
\textit{Security} & Istio, APISIX, Dex, oauth2-proxy, SpiceDB, OpenBao, ESO, Kyverno, Falco, Trivy, Velero, Chaos Mesh & \texttt{security}, \texttt{istio-system}, \texttt{external-secrets}, \texttt{velero} \\
\hline
\end{tabular}
\end{table}


\subsection{\textit{Data Ingestion and Processing Layer}}
\label{subsec:impl-ingestasi}

Apache Kafka \autocite{apachekafka2024} diterapkan menggunakan operator Strimzi. \textit{Kafka Connect} digunakan untuk integrasi dengan basis data eksternal melalui Debezium dan untuk \textit{sink} ke Apache Iceberg pada MinIO \autocite{minio2024} dengan katalog Lakekeeper \autocite{lakekeeper2025}. \textit{Karapace} berperan sebagai registri skema dengan dukungan Avro dan JSON Schema. Trafik masuk dari konsumen di luar \textit{mesh} dilewatkan melalui \textit{API gateway} APISIX \autocite{apisix2025} dengan mTLS yang dijembatani pada \textit{mesh edge}.

Konfigurasi \textit{cluster} Kafka dirangkum pada Tabel~\ref{tab:impl-kafka-config} agar pilihan teknis dan alasannya terbaca dalam satu pandangan. Keputusan pada tabel tersebut menjaga kontrak \textit{ingestion} tetap deklaratif, karena \texttt{KafkaTopic} dan \texttt{KafkaUser} CR membuat \textit{topic} beserta ACL dapat diaudit sebagai sumber daya Kubernetes, sedangkan kredensial SCRAM-SHA-512 mengalir dari OpenBao tanpa menyentuh manifes. \textit{Scaling} konsumen mengikuti panjang antrean nyata karena KEDA membaca \textit{lag} dari Kafka Exporter melalui Prometheus. Penempatan \textit{broker} di luar \textit{mesh} Istio dengan enkripsi SASL\_SSL \textit{built-in} Kafka mengikuti rancangan \textit{ingress flow} dari luar \textit{mesh} pada Subbab~\ref{subsec:lapisan-keamanan}, sehingga penyedia data tanpa \textit{sidecar} tetap terlayani tanpa melemahkan \textit{default} \textit{strict} internal.

% Konfigurasi jalur ingestion Kafka pada V.2.1 dalam bentuk tabel agar
% paragraf implementasi tidak menjadi tumpukan sintaks. Fakta dari
% platform/components/data-ingestion. Lebar total = 2.6+5.3+5.4 = 13.3cm
% + 6*tabcolsep(3pt) ~= 13.93cm (< \textwidth 14.0cm).
\begin{table}[!htb]
\centering
\caption{Konfigurasi Jalur Masuk Kafka pada Lingkungan Implementasi}
\label{tab:impl-kafka-config}
\footnotesize
\setlength{\tabcolsep}{3pt}
\begin{tabular}{|>{\raggedright\arraybackslash}m{2.6cm}|>{\raggedright\arraybackslash}m{5.3cm}|>{\raggedright\arraybackslash}m{5.4cm}|}
\hline
\centering\arraybackslash\textbf{Aspek} & \centering\arraybackslash\textbf{Konfigurasi} & \centering\arraybackslash\textbf{Alasan} \\
\hline
Mode \textit{cluster} & \textit{KRaft} murni tanpa ZooKeeper & Topologi tetap ringan pada satu \textit{node} \\
\hline
Protokol klien & SASL\_SSL dengan SCRAM-SHA-512, kredensial dari OpenBao melalui External Secrets Operator & Autentikasi dan enkripsi seragam tanpa rahasia pada manifes \\
\hline
Kontrak \textit{ingestion} & \texttt{KafkaTopic} dan \texttt{KafkaUser} CR beserta ACL & Kontrak tersimpan sebagai sumber daya Kubernetes yang dapat diaudit \\
\hline
Topologi \textit{topic} & Enam \textit{partition}, faktor replikasi satu, dinaikkan tiga pada multi-\textit{node} & Paralelisme konsumen tanpa melampaui kapasitas satu \textit{node} \\
\hline
Penskalaan konsumen & KEDA \texttt{ScaledObject} membaca \textit{lag} dari Kafka Exporter melalui Prometheus & Skala konsumen mengikuti panjang antrean nyata \\
\hline
Akses luar \textit{mesh} & \textit{Broker} di luar \textit{mesh} Istio (\texttt{sidecar.istio.io/inject: false}) dengan \textit{DestinationRule} menonaktifkan mTLS menuju \textit{bootstrap} & Produsen tanpa \textit{sidecar} tetap terlayani, enkripsi ditangani SASL\_SSL Kafka \\
\hline
\end{tabular}
\end{table}


\textit{Processing pipeline} pada \textit{layer} yang sama dirinci sebagai berikut. Pemrosesan \textit{stream} menggunakan operator Flink yang mengelola \textit{FlinkDeployment}. Pemrosesan \textit{batch} menggunakan operator Spark dan dijalankan sebagai \textit{SparkApplication}. Transformasi pada \textit{warehouse} menggunakan dbt yang dijalankan sebagai \textit{task} \texttt{KubernetesPodOperator} pada DAG Airflow \autocite{airflow2024}. Airflow mengorkestrasi \textit{pipeline} data berorientasi \textit{batch}, termasuk pembangunan \textit{gold layer}, sedangkan orkestrasi \textit{model training} ditangani oleh Kubeflow Pipelines.

Flink dipasang dengan \textit{session cluster} mode \textit{Standalone} dan tugas \textit{stream} dideklarasikan sebagai \texttt{FlinkSessionJob} CR sehingga setiap \textit{job} dapat di-\textit{deploy} secara independen melalui Argo CD. Spark Operator memanfaatkan \textit{webhook} mutasi untuk menyisipkan \textit{node selector} dan \textit{toleration} berdasarkan label \texttt{queue-name} yang ditetapkan oleh Kueue. dbt menjalankan model SQL pada ClickHouse dengan \texttt{profiles.yaml} yang merujuk \textit{secret} dari OpenBao melalui ESO. Airflow memakai \textit{git-sync} \textit{sidecar} yang menarik DAG dari Gitea sehingga pembaruan DAG tidak memerlukan pembangunan ulang \textit{image}. \textit{Checkpoint} dan \textit{savepoint} Flink ditulis ke MinIO sehingga pemulihan \textit{job} tidak kehilangan \textit{state}, dan \textit{topic} hasil \textit{stream} dikonsumsi \textit{analytics layer} melalui tabel \textit{Kafka engine} ClickHouse.

\subsection{\textit{Data Storage Layer}}
\label{subsec:impl-penyimpanan}

\textit{Warehouse} kolumnar menggunakan ClickHouse \autocite{clickhouse2024} dengan operator ClickHouse pada satu \textit{shard}. \textit{Data lake} menggunakan MinIO sebagai \textit{object store} dan Apache Iceberg sebagai \textit{table format}, dengan Lakekeeper \autocite{lakekeeper2025} sebagai katalog Iceberg REST dan Trino \autocite{trino2024} sebagai \textit{query engine}. Penyimpanan \textit{key-value} memakai Valkey \autocite{valkey2025} dengan protokol RESP untuk akses berlatensi rendah. Penyimpanan vektor memakai Qdrant \autocite{qdrant2025} sebagai basis data untuk \textit{embedding} berdimensi tinggi. Penyimpanan relasional memakai PostgreSQL melalui operator CNPG dengan \textit{plugin} \textit{barman-cloud} untuk pencadangan ke MinIO. MySQL dipasang sebagai \texttt{Deployment} dengan \textit{persistent volume} pada \textit{storage class} lokal untuk komponen yang terikat pada protokol MySQL. Apache Superset \autocite{superset2025} disiapkan sebagai antarmuka BI untuk peran \textit{business user} dengan koneksi ke ClickHouse, Trino, dan MySQL melalui SQLAlchemy.

ClickHouse memakai \texttt{ClickHouseKeeper} sebagai pengganti ZooKeeper untuk koordinasi metadata replikasi. MinIO dipasang dengan empat \texttt{drive} pada satu \textit{node} untuk menjaga skema penamaan \texttt{erasure code} yang seragam ketika \textit{cluster} berpindah ke banyak \textit{node}. Lakekeeper menyediakan \texttt{Warehouse} dan \texttt{Project} REST API yang dipakai oleh Spark, Trino, dan Flink dengan menggunakan \textit{header} \texttt{x-project-id} sebagai pemisah antar \textit{warehouse}. lakeFS \autocite{lakefs2025} dipasang sebagai \textit{version control layer} di atas MinIO sehingga eksperimen pada cabang data dapat dijalankan tanpa mengganggu data utama. MinIO disusun dengan \textit{bucket} terpisah per kepentingan sebagaimana dirangkum Tabel~\ref{tab:impl-minio-bucket}, dengan kapasitas yang dipantau melalui metrik MinIO pada Prometheus. Qdrant dilengkapi \textit{persistent volume}, \texttt{PodDisruptionBudget}, dan kunci API yang disalurkan External Secrets Operator.

% Tata letak bucket MinIO pada V.2.2 dalam bentuk tabel agar enumerasi
% bucket tidak memanjang di prosa. Fakta dari platform/components/storage.
% Lebar total = 3.4+10.12 = 13.52cm + 4*tabcolsep(3pt) ~= 13.94cm (< \textwidth).
\begin{table}[H]
\centering
\caption{Tata Letak \textit{Bucket} MinIO per Kepentingan}
\label{tab:impl-minio-bucket}
\footnotesize
\setlength{\tabcolsep}{3pt}
\begin{tabular}{|>{\raggedright\arraybackslash}m{3.4cm}|>{\raggedright\arraybackslash}m{10.12cm}|}
\hline
\textbf{\textit{Bucket}} & \textbf{Kepentingan} \\
\hline
\texttt{mlflow} & Artefak eksperimen dan model dari \textit{tracking} MLflow \\
\hline
\texttt{feast} & Registri fitur berbasis berkas (\texttt{registry.db}) yang dibagikan seluruh klien Feast \\
\hline
\texttt{warehouse} & Data \textit{lakehouse} berformat Iceberg di bawah katalog Lakekeeper \\
\hline
\texttt{flink-checkpoints} & \textit{Checkpoint} dan \textit{savepoint} \textit{job} \textit{streaming} Flink \\
\hline
\texttt{velero-backups} & Pencadangan terjadwal sumber daya \textit{cluster} dan \textit{persistent volume} \\
\hline
\end{tabular}
\end{table}


\subsection{\textit{Model Lifecycle Layer}}
\label{subsec:impl-model}

MLflow \autocite{mlflow2024} dijalankan dengan basis data PostgreSQL dan \textit{artifact store} pada MinIO. Kubeflow Pipelines \autocite{kubeflow2024} dijalankan pada \textit{namespace} terpisah dengan dukungan \textit{multi-tenant}, dilengkapi Kubeflow Notebooks untuk lingkungan eksperimen interaktif dan Kubeflow Trainer \autocite{kubeflowtrainer2025} untuk \textit{distributed training}. KServe \autocite{kserve2024} dipakai untuk \textit{inference} dengan dukungan \textit{runtime} MLflow, scikit-learn, XGBoost, PyTorch, dan ONNX, dan dengan Knative Serving \autocite{knative2025} sebagai \textit{serverless layer} \textit{built-in} yang menyediakan \textit{scale-to-zero}. Argo Workflows berperan sebagai \textit{orchestrator} \textit{job} di luar Kubeflow Pipelines, sedangkan Argo Rollouts \autocite{argorollouts2025} mengelola promosi versi pada \textit{serving layer} dengan strategi \textit{canary} dan \textit{progressive analysis} berbasis metrik Prometheus, sehingga otomatisasi \textit{deployment} model pada KF-06 terpenuhi.

Kubeflow Pipelines memakai MySQL sebagai \textit{metadata store} dan MinIO sebagai \textit{artifact store}, dengan \texttt{kfp-launcher} yang membungkus eksekusi langkah \textit{pipeline} agar autentikasi ke MinIO dilakukan menggunakan \texttt{secretAsEnv}. KServe \texttt{InferenceService} dideklarasikan pada \textit{namespace} \texttt{model-serving} dengan label \texttt{part-of=kserve} agar Kyverno tidak menolak \textit{Pod} \textit{mutator} \textit{built-in}. Kueue \autocite{kueue2024} mengelola antrean \textit{training} dengan \texttt{ClusterQueue} per kelompok \textit{workload} sehingga sumber daya CPU dan GPU dapat dibagi adil antara \textit{tenant}, memenuhi KF-14 tentang manajemen sumber daya. Argo Rollouts mengevaluasi metrik dari Prometheus pada setiap langkah \textit{canary} untuk menentukan apakah versi \textit{canary} layak dipromosikan menjadi versi \textit{stable}.

\subsection{\textit{Observability Layer}}
\label{subsec:impl-observability}

Prometheus \autocite{prometheus2024} mengumpulkan metrik dari semua komponen dengan \textit{ServiceMonitor} dan \textit{PodMonitor}. Loki \autocite{loki2024} mengumpulkan log melalui Alloy. Grafana Tempo \autocite{tempo2024} menerima \textit{trace} dari OpenTelemetry Collector yang tersebar pada \textit{node}. Grafana \autocite{grafana2024} menggabungkan ketiganya pada satu antarmuka.

Sloth \autocite{sloth2025} menerjemahkan definisi \textit{service-level objective} ke dalam aturan Prometheus sehingga \textit{error budget} dapat dipantau bersama metrik lain. Pushgateway berperan sebagai komponen pengumpul metrik untuk \textit{job} berdurasi pendek yang berjalan di luar siklus \textit{scrape}, termasuk CronJob pengukuran \textit{drift} dan \textit{CronWorkflow} pemicu \textit{retraining}. OpenCost dipakai untuk pemantauan biaya per \textit{workload}, sedangkan Pyroscope mengumpulkan profil CPU dan memori secara berkelanjutan dari layanan yang menjadi \textit{hotspot}. Evidently \autocite{evidently2025} dijalankan sebagai \textit{job} berkala yang menghasilkan \textit{report} HTML kualitas distribusi fitur dan disimpan pada MinIO sebagai jejak audit yang dapat dirujuk dari DataHub. Pemenuhan kebutuhan yang bergantung pada arsitektur terintegrasi ini diverifikasi melalui skenario evaluasi pada Bab~\ref{chap:evaluasi}.

\section{Implementasi Sub-sistem Tata Kelola Data}
\label{sec:impl-tata-kelola}

Subbab ini menurunkan rancangan tata kelola pada Bagian~\ref{sec:tata-kelola} menjadi implementasi sebagai pemenuhan T-2, dengan fokus pada konsistensi katalog metadata, jadwal pengujian kontrak data, dan penegakan kontrol akses pada tiga \textit{layer} yang berbeda. Penomoran \texttt{KF-08}, \texttt{KF-09}, dan \texttt{KF-10}, serta kebutuhan KNF terkait keamanan dipakai sebagai acuan untuk menilai apakah implementasi memenuhi seluruh kebutuhan yang dirumuskan pada Bab~\ref{chap:analisis-masalah}. Katalog diwujudkan oleh DataHub dengan penerimaan \textit{lineage} OpenLineage, kontrak kualitas oleh Great Expectations pada \textit{pipeline} Airflow, dan kontrol akses oleh gabungan dex, SpiceDB, serta kebijakan jaringan. Autentikasi konsol, otorisasi granular, dan kebijakan jaringan antar \textit{namespace} membentuk ketiga \textit{access layer} yang dimaksud.

\subsection{Katalog Metadata dan \textit{Lineage}}
\label{subsec:impl-katalog}

DataHub \autocite{datahub2024} dijalankan dengan komponen GMS, \textit{frontend}, dan beberapa \textit{job} \textit{ingest}, dengan OpenSearch sebagai \textit{search backend} dan PostgreSQL sebagai \textit{metadata backend}. \textit{Job ingest} berjalan secara berkala sebagai \textit{CronJob} untuk menarik metadata dari Kafka, ClickHouse, MinIO dan Iceberg, Feast, MLflow, dan Kubeflow Pipelines. Hasil \textit{ingest} menyatukan \textit{lineage} antar \textit{layer}, mulai dari \textit{topic} pada Kafka hingga model yang dilayani oleh KServe. OpenLineage \autocite{openlineage2025} dipakai sebagai protokol pengirim \textit{event} \textit{lineage} dari komponen \textit{pipeline}, baik Airflow maupun Kubeflow Pipelines, agar graf ketertelusuran di DataHub tersusun dengan kosakata yang seragam.

\textit{Image} \texttt{datahub-feast-deps} dibangun secara lokal dengan Kaniko untuk menambahkan \textit{plugin} Feast pada \textit{job ingest} sehingga sumber Feast dapat dipindai secara langsung. Autentikasi ke GMS menggunakan token sistem yang disimpan di OpenBao dan disebarkan ke \textit{namespace} \texttt{data-governance} oleh \textit{External Secrets Operator}. Konfigurasi DataHub pada level UI dikelola melalui \texttt{datahub-frontend-secret} dengan \textit{checksum annotation} pada GMS dan \textit{frontend} agar rotasi \textit{secret} memicu \textit{auto-roll} \textit{deployment}. Jadwal \textit{ingest} dapat berkonflik dengan \textit{workload} besar di \textit{node}, dan untuk itu \textit{native sidecar} Istio diaktifkan pada \textit{CronJob} agar koneksi mTLS ke GMS tidak putus saat \textit{sidecar} dihentikan setelah \textit{job} selesai. \textit{Job} berdurasi pendek lain yang berjalan di luar \textit{mesh} mengirim metriknya melalui titik \texttt{PERMISSIVE} Pushgateway sesuai rancangan pengecualian pada Subbab~\ref{subsec:lapisan-keamanan}, sehingga hanya \textit{job} yang benar-benar membutuhkan koneksi mTLS ke layanan dalam \textit{mesh} yang membawa \textit{native sidecar}.

\subsection{Kontrak Data dan Kualitas}
\label{subsec:impl-kualitas}

Great Expectations \autocite{greatexpectations2024} mendefinisikan ekspektasi kualitas data pada \textit{bronze layer} setelah \textit{ingestion} dan pada \textit{silver} dan \textit{gold layer} setelah transformasi. Setiap \textit{expectation suite} disimpan pada repositori Git dan dijalankan oleh \textit{CronJob} yang membaca hasil ke DataHub. Pengujian yang gagal akan menutup \textit{promotion flow} ke \textit{layer} berikutnya dan memunculkan peringatan pada Grafana. Hasil ekspektasi disimpan ke MinIO sebagai \textit{Data Docs} HTML sehingga konteks kegagalan dapat ditelusuri tanpa perlu menjalankan kembali \textit{suite}.

dbt menjalankan model SQL untuk \textit{silver} dan \textit{gold layer} pada ClickHouse, dengan \texttt{schema.yml} yang memuat \textit{tests} dasar seperti \texttt{unique}, \texttt{not\_null}, dan \texttt{accepted\_values}. Hasil dbt dirilis ke DataHub melalui \textit{run event} OpenLineage yang dikirim oleh \textit{task} \texttt{KubernetesPodOperator} dbt pada DAG Airflow. \textit{Checkpoint} Great Expectations dijadwalkan per \textit{layer} dengan jeda yang berbeda agar ekspektasi yang berat tidak berdampak pada beban \textit{warehouse}. \textit{Suite} fakta \textit{training} diuji setelah materialisasi pada tabel \texttt{gold.fct\_training\_data} agar \textit{training set} tidak masuk \textit{training pipeline} apabila kualitas data tidak memenuhi kontrak.

\subsection{Kontrol Akses}
\label{subsec:impl-akses}

Kontrol akses pada \textit{network layer} menggunakan \textit{NetworkPolicy} dengan kebijakan \textit{default-deny} per \textit{namespace} dan pengecualian yang dideklarasikan secara eksplisit. Pada \textit{application layer}, kontrol akses memakai \textit{authentication flow} tunggal dex dan oauth2-proxy yang telah dirinci pada Subbab~\ref{subsec:rahasia-identitas}, sehingga akses ke konsol tata kelola tunduk pada sesi identitas yang sama dengan layanan platform lain. Pada \textit{admission layer}, Kyverno \autocite{kyverno2024} menegakkan kebijakan keamanan \textit{workload} sesuai rancangan pada Subbab~\ref{subsec:identitas-akses}. Falco \autocite{falco2025} memantau \textit{event} \textit{runtime} pada level \textit{kernel} untuk mendeteksi pelanggaran kebijakan keamanan setelah \textit{pod} berjalan.

SpiceDB \autocite{spicedb2025} berfungsi sebagai layanan otorisasi yang dipakai DataHub dan Kubeflow Pipelines untuk hubungan akses pada level objek, dengan basis data PostgreSQL melalui CNPG. \textit{Policy} Kyverno disusun pada \textit{namespace} \texttt{security} dengan \textit{policy generator} yang menyebarkan kebijakan ke \textit{namespace} aplikasi tanpa duplikasi sumber daya. Kyverno mode \texttt{failurePolicy=Ignore} dipakai pada \texttt{ValidatingPolicy} agar admisi tidak terblokir saat \textit{webhook} gagal pada kondisi \textit{cluster} yang sibuk. Falco mengirim peringatan ke Loki melalui \textit{falcosidekick} dengan \textit{topic} khusus sehingga tim keamanan dapat menelusuri \textit{event} terkait kebijakan dari satu antarmuka Grafana. Pemenuhan KF-08, KF-09, dan KF-10 diverifikasi melalui skenario evaluasi pada Bab~\ref{chap:evaluasi}.

\section{Implementasi Sub-sistem Deteksi \textit{Drift} dan \textit{Continuous Training}}
\label{sec:impl-drift}

Implementasi sub-sistem deteksi \textit{drift} menindaklanjuti rancangan pada Bagian~\ref{sec:drift-ct} sebagai pemenuhan T-3, dengan fokus pada otomatisasi penuh \textit{control loop} dari pengukuran sampai promosi versi serta penegakan ambang yang dapat dikonfigurasi per model. Pemicu \textit{retraining} dirancang agar dapat dijalankan secara berkala atau dipicu oleh \textit{event} \textit{drift} yang melebihi ambang. Keseluruhan implementasi pada bagian ini menjawab KF-07 untuk deteksi \textit{drift} otomatis berikut \textit{trigger} \textit{retraining}, sehingga model pada \textit{serving layer} tidak dibiarkan usang tanpa evaluasi berkala saat distribusi data bergeser. Seluruh langkah \textit{loop} diwujudkan pada \texttt{CronWorkflow} \texttt{retrain-on-drift} yang membaca skor dari tabel analitik dan memanggil \textit{training pipeline} Kubeflow Pipelines.

\subsection{Detektor \textit{Drift}}
\label{subsec:impl-detektor}

Detektor \textit{drift} diimplementasikan sebagai \textit{service} Python yang membaca fitur dari ClickHouse dengan \textit{time window} yang dapat dikonfigurasi. Indikator PSI dihitung dengan \textit{binning} adaptif, dan uji Kolmogorov-Smirnov dijalankan dengan ambang yang dapat disesuaikan per fitur \autocite{reis2016kdd, bayram2022concept, lu2019learning}. Hasil pengukuran ditulis ke ClickHouse sebagai tabel \texttt{gold.drift\_multi\_scale} dan dipublikasikan ke Pushgateway agar dapat diserap oleh Prometheus. Evidently \autocite{evidently2025} dijalankan sebagai pelapor pelengkap yang menghasilkan ringkasan distribusi fitur dan kinerja model dalam bentuk \textit{report} HTML yang dapat diaudit dan dilampirkan pada katalog DataHub.

\textit{Service} berjalan sebagai \textit{CronJob} pada \textit{namespace} \texttt{model-lifecycle} dengan jadwal lima belas menit untuk fitur fluktuatif dan satu jam untuk \textit{tabular feature} yang lebih stabil. Konfigurasi ambang per fitur dideklarasikan pada \texttt{ConfigMap} \texttt{drift-thresholds} sehingga perubahan ambang dapat dilakukan melalui \textit{commit} pada Git tanpa membangun ulang \textit{image}. Hasil tiap eksekusi ditulis dengan \texttt{checksum} versi konfigurasi sehingga riwayat ambang dapat ditelusuri saat audit. \textit{Image} detektor dibangun dengan \texttt{uv} pada Dockerfile multi-stage agar model SBERT untuk deteksi \textit{drift} \textit{embedding} disertakan langsung pada \textit{layer} terakhir \textit{image}, sehingga \textit{cold start} tidak memerlukan unduhan model.

\subsection{\textit{Retraining Control Loop}}
\label{subsec:impl-control-loop}

\textit{Control loop} dijalankan oleh Argo CronWorkflow dengan jadwal yang dapat dikonfigurasi per model. \textit{Workflow} menjalankan dua langkah utama, yaitu pengukuran \textit{drift} dan pemicu \textit{training pipeline}. Pemicu \textit{training} mengirim permintaan ke Kubeflow Pipelines sebagai pemenuhan KF-05 tentang otomatisasi \textit{training pipeline}. \textit{Training pipeline} membaca fitur \textit{training} dari tabel \textit{gold} ClickHouse, menjalankan \textit{training} pada \textit{job} dengan sumber daya yang dialokasikan oleh Kueue \autocite{kueue2024}, lalu mencatat metrik dan artefak model ke MLflow. Kelayakan model baru untuk menggantikan versi produksi ditentukan pada tahap promosi \textit{canary} yang menilai metrik \textit{serving}-nya secara langsung sebelum trafik penuh dialihkan, bukan pada perbandingan tersendiri sebelum \textit{deployment}.

Pada tahap promosi, Argo Rollouts \autocite{argorollouts2025} menyalurkan trafik \textit{serving layer} secara bertahap ke versi baru memakai strategi \textit{canary} dengan analisis metrik Prometheus, sehingga \textit{rollback} otomatis dapat dilakukan ketika kinerja menurun. \texttt{retrain-on-drift} CronWorkflow disimpan pada direktori \textit{workflow} milik \textit{use case} yang terpisah dari \texttt{platform/} agar dapat di-\textit{deploy} bersama \textit{use case} ketika pengujian dijalankan. Konfigurasi penonaktifan \textit{cache} KFP diterapkan pada langkah \texttt{train} dan \texttt{deploy} agar setiap eksekusi memakai data terbaru tanpa dipengaruhi \textit{cache} historis. Pengiriman metrik \textit{workflow} dilakukan di dalam langkah \texttt{decide-and-trigger} yang menulis ke Pushgateway, sehingga seluruh metrik eksekusi tetap terkumpul pada Prometheus meskipun \textit{workflow} berdurasi singkat.

\subsection{Penanganan Kasus Khusus}
\label{subsec:impl-kasus-khusus}

Implementasi menangani tiga kasus khusus yang sama dengan rancangan pada Subbab~\ref{subsec:kasus-khusus}. Perwujudan ketiganya tidak menambah komponen baru di luar \texttt{CronWorkflow} dan Argo Rollouts yang sudah terpasang. Perilaku tiap kasus dapat diperiksa dari manifes dan log eksekusi pada \textit{namespace} \texttt{model-lifecycle}. Rinciannya sebagai berikut.

\vspace{0.5em}
\begin{enumerate}
    \item \textbf{Waktu tunggu label per model}

    Waktu tunggu label diatur per model agar \textit{retraining} tidak dipicu sebelum label tersedia cukup.

    \vspace{0.5em}

    \item \textbf{\textit{Sliding window} untuk meredam fluktuasi}

    Ambang \textit{drift} dihitung dengan \textit{sliding window} untuk meredam fluktuasi singkat.

    \vspace{0.5em}

    \item \textbf{Penghentian promosi saat analisis \textit{canary} gagal}

    Ketika analisis \textit{canary} gagal, promosi dihentikan sehingga versi stabil pada \textit{serving layer} tetap dipertahankan.
\end{enumerate}

Penanganan ketiga kasus khusus mengikuti rancangan pada Subbab~\ref{subsec:kasus-khusus}. Ambang PSI dan KS beserta \textit{lookback window} mengikuti mekanisme \texttt{ConfigMap} \texttt{drift-thresholds} yang telah dirinci pada Subbab~\ref{subsec:impl-detektor}, termasuk jejak \texttt{checksum} versi konfigurasinya. Nilai parameter yang aktif dapat dibaca langsung dari manifes sehingga audit kebijakan tidak memerlukan akses ke sistem lain. Pemenuhan KF-07 diverifikasi melalui skenario evaluasi pada Bab~\ref{chap:evaluasi}.

\section{Implementasi Sub-sistem Layanan Fitur \textit{Dual-Store}}
\label{sec:impl-feast}

Sub-sistem layanan fitur \textit{dual-store} mewujudkan rancangan pada Bagian~\ref{sec:feature-store-dual} sebagai pemenuhan T-4, dengan satu sumber definisi fitur dan satu kontrak Feast untuk \textit{tabular feature}, sedangkan fitur \textit{embedding} dikelola pada indeks vektor terpisah. \textit{Tabular feature} dimaterialisasi ke \textit{online store} Valkey dan diakses melalui API Feast, sedangkan indeks vektor Qdrant diisi dan diakses melalui klien \textit{native} di luar materialisasi Feast. Sumber definisi tunggal di \texttt{feature\_repo/} mengikat ketiga penyimpanan itu sehingga perubahan skema fitur cukup dilakukan pada satu tempat. Bagian ini menjawab KF-01 sampai KF-04 yang menuntut manajemen fitur terpusat, \textit{serving} ganda \textit{online}-\textit{offline}, validitas temporal, dan dukungan \textit{vector embedding}.

\subsection{Konfigurasi Feast}
\label{subsec:impl-feast-config}

Feast \autocite{feast2024} dikonfigurasi dengan \texttt{feature\_store.yaml} yang menentukan \textit{offline store} ClickHouse, \textit{online store} Valkey \autocite{valkey2025} melalui protokol RESP, dan registri \textit{file-based} pada MinIO. Definisi fitur ditulis dalam Python pada modul \texttt{feature\_definitions.py} dan didaftarkan ke registri melalui \texttt{feast apply}. Materialisasi fitur ke \textit{online store} dijalankan oleh \textit{job} berkala yang membaca dari ClickHouse dengan rentang waktu yang sesuai. Berkas \texttt{feature\_store.yaml} yang sama dipakai oleh seluruh klien sehingga tidak ada konfigurasi penyimpanan yang didefinisikan ulang per layanan.

Registri Feast disimpan sebagai \texttt{registry.db} di \textit{bucket} MinIO dengan kunci yang dirotasi melalui ESO sehingga klien Feast pada berbagai \textit{namespace} berbagi sumber rujukan yang sama. \textit{Service} \texttt{feast serve} dipasang sebagai \textit{Deployment} yang dijangkau melalui \texttt{Service} \texttt{feast-server} pada \textit{namespace} \texttt{model-lifecycle}. Pembaruan definisi fitur dilakukan oleh \textit{service} materialisasi yang menjalankan \texttt{feast apply} sebagai langkah pada DAG \texttt{data\_pipeline} Airflow, sehingga registri selalu sinkron dengan definisi pada \texttt{feature\_repo/}. \textit{Client} Feast pada layanan pemanggil dikonfigurasi dengan \texttt{registry\_path} URI \texttt{s3://}, sehingga \textit{single source of truth} tetap terjaga ketika klien berjalan pada \textit{namespace} berbeda. Layanan pemanggil memakai SDK Python yang seragam, yaitu klien Feast untuk \textit{online feature}, klien MLflow untuk registri model, dan protokol KServe v2 untuk \textit{inference}, dengan \textit{environment} \texttt{uv} yang mengunci versi ketiganya pada tiap layanan, sehingga KF-13 tentang API dan SDK terpadu terpenuhi.

\subsection{Dukungan \textit{Vector Embeddings}}
\label{subsec:impl-vektor}

\textit{Vector embedding} disimpan sebagai fitur dengan tipe \texttt{Array(Float32)} pada ClickHouse dan sebagai \textit{point} berdimensi 768 pada \textit{collection} Qdrant \autocite{qdrant2025}. Pengindeksan HNSW diaktifkan pada Qdrant dengan parameter \texttt{m} dan \texttt{ef\_construct} yang dapat dikonfigurasi, dan kueri tetangga terdekat dilayani melalui API gRPC dengan filter \textit{payload} pada saat pencarian. \textit{Service} \textit{vector-embedding} pada \texttt{platform/services/processing/vector/} membungkus model \texttt{sentence-transformers/all-mpnet-base-v2} sebagai penghasil \textit{embedding} dan menyajikan API yang dipanggil oleh \textit{stream processor} pada saat \textit{event} masuk \autocite{reimers2019sentence}. Metrik jarak yang dipakai adalah kosinus sesuai \textit{default} layanan \textit{embedding} untuk vektor yang dinormalkan.

\textit{Image} \textit{vector-embedding} dibangun dengan pola \textit{bake} model yang sama seperti \textit{image} detektor \textit{drift} pada Subbab~\ref{subsec:impl-detektor}, sehingga \textit{startup probe} tidak terganggu oleh unduhan model pada pemicu pertama. \texttt{Collection} Qdrant dideklarasikan oleh \textit{init Job} yang mengirim \texttt{PUT} ke \texttt{/collections/}\textit{nama-koleksi} dengan parameter HNSW sesuai konfigurasi. \textit{Snapshot} koleksi disimpan ke MinIO secara berkala dengan \textit{Job} \texttt{qdrant-snapshot} yang menambah objek pencadangan ke jadwal Velero. Koleksi vektor Qdrant diakses langsung oleh layanan \textit{embedding} melalui klien \textit{native} Qdrant, terpisah dari \textit{online store} Valkey yang menyajikan \textit{tabular feature} melalui Feast, sehingga \textit{vector query flow} dan \textit{tabular feature flow} tidak saling bergantung.

\subsection{Materialisasi dan \textit{Point-in-Time Correctness}}
\label{subsec:impl-pit}

\texttt{get\_historical\_features} dijalankan pada \textit{training pipeline} dengan kolom \texttt{event\_timestamp} pada \textit{entity dataframe}. Feast melakukan \textit{point-in-time join} pada \textit{offline store} ClickHouse dengan memilih nilai fitur yang \texttt{timestamp}-nya tidak melampaui \texttt{event\_timestamp} setiap baris entitas dan masih berada dalam rentang \texttt{ttl} \textit{FeatureView}. Ketika terdapat lebih dari satu kandidat untuk satu titik waktu, kolom \texttt{created\_timestamp} dipakai sebagai penentu agar hanya nilai paling mutakhir yang diambil, sehingga nilai fitur yang dipakai pada saat \textit{training} identik dengan nilai yang tersedia pada saat itu. Materialisasi ke \textit{online store} berjalan secara berkala dengan \texttt{feast materialize-incremental}, dan \textit{stream materializer} pada Flink menulis fitur \textit{real-time} ke Valkey untuk \textit{tabular feature} secara bersamaan dengan penulisan ke ClickHouse. Keterpaduan \textit{batch pipeline} dan \textit{stream} pada satu kontrak Feast ini memenuhi KF-11.

Materialisasi dijalankan berkala oleh DAG Airflow melalui \textit{task} \texttt{feast\_materialize} yang mendorong fitur dari \textit{offline store} ClickHouse ke \textit{online store} Valkey pada jadwal per jam. Karena nilai pada Valkey merupakan hasil materialisasi dari ClickHouse, konsistensi antar \textit{store} terjaga secara struktural tanpa salinan yang dipelihara terpisah. Kualitas data diverifikasi oleh \textit{CronJob} \texttt{ge-checkpoint} yang menjalankan \textit{suite} Great Expectations terhadap tabel pada ClickHouse setiap enam jam, dengan hasil validasi ditulis ke MinIO dan \textit{gauge} lulus atau gagal dikirim ke Pushgateway. Pengujian \textit{point-in-time} dijalankan sebagai skenario SK-F-03 yang mengaudit \texttt{get\_historical\_features} pada \textit{training dataset}, memastikan kontrak \texttt{event\_timestamp} dan \texttt{ttl} tetap dipatuhi oleh \texttt{FeatureView} yang ditambahkan.

\section{Verifikasi Implementasi}
\label{sec:verifikasi}

Verifikasi \textit{execution flow} dilakukan dengan satu \textit{use case} domain nyata sebagai kasus pengujian, yaitu analitik pasar aset kripto, yang struktur layanan, konfigurasi, dan datanya dirinci pada Bab~\ref{chap:evaluasi}. Pemilihan domain ini berpijak pada tiga alasan yang dapat dirunut ke bukti. Alasan pertama terletak pada \textit{use case} ini yang menjalankan seluruh sembilan \textit{layer} representatif yang ditetapkan pada Subbab~\ref{sec:layanan-industri}, dengan pemetaan tiap \textit{layer} ke komponennya dirinci pada Subbab~\ref{subsec:instans-verifikasi}, sehingga verifikasi fungsionalnya menyentuh himpunan layanan yang sama dengan yang dituntut praktik industri. Alasan kedua muncul karena pasar aset kripto beroperasi terus menerus tanpa jam tutup bursa \autocite{makarov2020trading}, sehingga \textit{ingestion pipeline} dan pemrosesan \textit{stream} teruji oleh aliran \textit{event} nyata selama 24 jam penuh, bukan oleh beban sintetis yang dijadwalkan. Alasan ketiga bertumpu pada volatilitas harga aset kripto yang tinggi dan bergerombol \autocite{katsiampa2017volatility} yang membuat distribusi fitur turunan harga rawan bergeser dalam hitungan jam, sehingga sub-sistem deteksi \textit{drift} dan \textit{retraining} otomatis memperoleh pemicu alami tanpa perlu injeksi buatan sebagai satu-satunya sumber pengujian. Pemilihan ini tidak mengubah sifat domain-agnostik platform, karena seluruh kode dan konfigurasi platform tetap berada pada repositori \texttt{platform/}, sedangkan kode dan konfigurasi spesifik domain tinggal pada direktori \textit{use case} tersendiri yang mengikuti pola \textit{base} dan \textit{overlays} \textit{Kustomize}. Dengan pemisahan itu, \textit{use case} lain cukup menyalin dan menyesuaikan berkas konfigurasi domain tanpa mengubah kode platform, sehingga pemilihan kripto tidak mengurangi keberlakuan platform pada domain lain.

\subsection{Lingkup Verifikasi}
\label{subsec:lingkup-verifikasi}

Verifikasi mencakup lima \textit{flow} yang masing-masing menyasar satu sub-sistem. Kelima \textit{flow} itu dipilih dari alur ujung ke ujung \textit{use case} sehingga setiap kontrak antar \textit{layer} teruji oleh trafik nyata. Setiap \textit{flow} memiliki keluaran yang dapat diamati pada \textit{dashboard} atau tabel analitik sebagai bukti. Uraian kelima \textit{flow} tersebut adalah sebagai berikut.

\vspace{0.5em}
\begin{enumerate}
    \item \textbf{\textit{Ingestion pipeline}}

    \textit{Ingestion pipeline} diuji dengan \textit{stream} dari sumber eksternal yang masuk ke Kafka melalui layanan \texttt{websocket-collector} pada \textit{use case}.

    \vspace{0.5em}

    \item \textbf{\textit{Processing pipeline}}

    \textit{Processing pipeline} diuji dengan \texttt{stream-processor} berbasis Flink yang menulis fitur agregasi ke ClickHouse dan Valkey.

    \vspace{0.5em}

    \item \textbf{\textit{Training flow}}

    \textit{Training flow} diuji dengan \textit{pipeline} Kubeflow yang membaca fitur dari Feast dan menulis model ke MLflow.

    \vspace{0.5em}

    \item \textbf{\textit{Serving flow}}

    \textit{Serving flow} diuji dengan \textit{InferenceService} KServe yang memuat model dari MLflow, dengan promosi versi pada \textit{serving layer} yang dikelola oleh Argo Rollouts.

    \vspace{0.5em}

    \item \textbf{\textit{Retraining pipeline}}

    \textit{Retraining pipeline} diuji dengan menjalankan \textit{drift-detector} dan memicu \textit{retrain-on-drift} \textit{workflow}.
\end{enumerate}

Uji ketahanan pendek menggunakan Chaos Mesh \autocite{chaosmesh2025} untuk menginjeksikan \textit{fault} pada \textit{pod} terpilih. Skenario yang dipakai adalah \texttt{PodChaos} yang menghentikan paksa \textit{pod} terpilih lalu mengamati pemulihannya oleh \textit{controller}. Pencadangan dijalankan terjadwal oleh Velero \autocite{velero2025} ke \textit{bucket} MinIO khusus pencadangan sebagai langkah disiplin pemulihan. Pada skenario ketahanan di Bab~\ref{chap:evaluasi}, kedua mekanisme itu diuji kembali.

Pemilihan kelima \textit{flow} tersebut bertujuan agar setiap sub-sistem mendapat verifikasi minimal satu skenario nyata, sehingga ketidaktepatan rancangan dapat ditemukan sebelum penilaian akhir pada Bab~\ref{chap:evaluasi}. Skenario \textit{use case} menetapkan beban yang lebih tinggi pada \textit{ingestion pipeline} dan pemrosesan dibandingkan \textit{training flow}, karena \textit{use case} ber-\textit{throughput} tinggi memproduksi banyak \textit{event} per detik dari beberapa entitas secara paralel. Pengamatan dilakukan dengan \textit{dashboard} Grafana yang menampilkan metrik latensi, \textit{throughput}, panjang antrean Kafka, dan tingkat keberhasilan materialisasi. Hasil pengamatan dibandingkan dengan target metrik pada KNF Bab~\ref{chap:analisis-masalah} untuk menentukan apakah setiap \textit{flow} memenuhi target yang telah ditetapkan.

\subsection{Pengamatan Hasil Awal}
\label{subsec:pengamatan-awal}

Hasil pengamatan awal memperlihatkan bahwa \textit{control plane} platform terbentuk sesuai rancangan. Seluruh komponen pada kelima \textit{flow} berhasil di-\textit{deploy} dan direkonsiliasi oleh Argo CD, \textit{retraining control loop} terpicu otomatis ketika injeksi \textit{drift} terkendali diterapkan pada salah satu fitur, dan promosi versi pada \textit{serving layer} dikendalikan oleh Argo Rollouts sesuai konfigurasi \textit{canary}. Reproduksibilitas \textit{flow} diuji dengan merekonstruksi seluruh komponen platform dari Argo CD pada \textit{cluster} k3s yang baru, dan rekonstruksi tersebut berhasil dilakukan tanpa intervensi manual selain \textit{bootstrap} Argo CD awal. Penilaian kebenaran data \textit{end-to-end}, validitas \textit{point-in-time} pada \textit{training}, latensi \textit{serving}, \textit{throughput} pemrosesan, dan ketepatan deteksi \textit{drift} diuraikan secara terperinci pada Bab~\ref{chap:evaluasi}.

Pengamatan tambahan menunjukkan bahwa \textit{observability flow} berhasil mempertahankan metrik dan log dari seluruh komponen pada satu antarmuka Grafana, dan bahwa \textit{event} \textit{lineage} dari \textit{emitter} OpenLineage Airflow dan Spark serta konektor \textit{ingestion} DataHub untuk Feast, MLflow, ClickHouse, dan Kafka muncul pada satu graf metadata dengan kosakata yang seragam. Pencadangan Velero berjalan terjadwal tanpa intervensi, dan pengujian pemulihan dari pencadangan memperlihatkan bahwa konfigurasi seluruh komponen dapat direkonstruksi tanpa langkah manual. Uji \textit{fault injection} Chaos Mesh menunjukkan bahwa \textit{pod} yang dimatikan secara paksa dijadwalkan ulang secara otomatis oleh \textit{controller} Kubernetes, dan bahwa Argo Rollouts otomatis menahan promosi ketika metrik \textit{canary} menyentuh ambang. Catatan pengamatan ini menjadi dasar bagi evaluasi pada Bab~\ref{chap:evaluasi} yang menilai pemenuhan KF dan KNF terhadap target metrik.
